\section{Conclusion}
\label{sec:conclusion}

Flood risk follows valleys, not grids.
Standard machine-learning models for flood susceptibility ignore this ---
treating every location as independent, as if water does not flow downhill through
connected river valleys.
This study shows that encoding watershed connectivity changes the answer,
and that the improvement is large enough to matter for real infrastructure decisions.

We presented the first state-wide flash flood susceptibility assessment for
Himachal Pradesh that simultaneously addresses three major shortcomings of
existing work: single-basin coverage, inflated random-split validation,
and point-estimate risk maps with no uncertainty.

The principal findings are as follows:

\begin{enumerate}

  \item \textbf{Graph structure improves susceptibility prediction.}
    The GraphSAGE GNN, trained on a directed watershed connectivity graph,
    outperformed all pixel-based baselines (RF, XGBoost, LightGBM, stacking ensemble)
    in leave-one-basin-out spatial block cross-validation,
    with the largest gains in basin folds with long, structured upstream networks.
    This confirms the hypothesis that upstream--downstream flood propagation,
    invisible to standard ML models, is an independent predictor of
    local flash flood susceptibility.

  \item \textbf{Conformal prediction intervals reveal regime-dependent calibration.}
    The 90\% prediction intervals achieved 82.9\% empirical coverage
    on the held-out 2023 temporal test set.
    Coverage is high in low-susceptibility areas (96.3\%) but degrades
    for high-susceptibility locations (45--59\%), indicating that
    SAR label noise in the high-risk regime reduces conformal calibration quality.
    Despite undercoverage, the uncertainty width map remains operationally useful:
    very high susceptibility combined with narrow uncertainty identifies
    the 4,409\,km$^2$ highest-priority zones for pre-monsoon intervention.

  \item \textbf{Spatial cross-validation corrects AUC inflation.}
    Spatial block CV (leave-one-block-out) yields pixel-based model AUC of
    0.870--0.881, versus the \citet{saha2023} benchmark of 0.88 with
    random splits --- a modest but more honestly obtained comparison.
    The benchmark's small AUC advantage likely reflects the 5--15\% inflation
    documented for random splits with spatially autocorrelated data
    \citep{valavi2019, roberts2017cross}, rather than a genuinely superior model.

  \item \textbf{Elevation, plan curvature, and slope are the dominant drivers.}
    SHAP analysis identified these three terrain morphometry factors as
    accounting for 73\% of total predictor importance across HP,
    with marked district-level variation.
    Rainfall and TWI rank 4th--5th, indicating that topographic position
    is the primary filter for flash flood susceptibility in HP's mountainous terrain.

  \item \textbf{High susceptibility covers 14.3\% of HP's study domain.}
    High and Very High susceptibility zones total 15,785\,km$^2$,
    concentrated in the lower Beas, middle Satluj, and Ravi valley corridors.
    OSM-based infrastructure exposure analysis identifies 1,457\,km of highways,
    2,759 bridges, and 92 tourist accommodation units within High or Very High zones,
    with 40 settlements in the Very High category.

\end{enumerate}

\subsection*{Policy Implications for HP State Disaster Management Authority}

The primary operational output of this study is a three-layer geospatial product:
(1) a point-estimate susceptibility map classified into four levels;
(2) an uncertainty width map indicating confidence in each classification;
and (3) an infrastructure exposure overlay.
Together, these layers enable HP SDMA to:
(a) identify the 5--10 highest-priority corridors for pre-monsoon clearance
    and early warning system installation;
(b) prioritise capital repair spending on bridges and roads in
    narrow-uncertainty, high-susceptibility zones; and
(c) issue targeted tourism advisories for the July--August peak season,
    distinguishing high-certainty risk zones (immediate action needed)
    from high-uncertainty risk zones (precautionary monitoring needed).

\subsection*{Directions for Future Research}

Three directions offer the highest marginal return:
(1) extending the SAR inventory to cover the full Sentinel-1 archive (2016--2024)
and adding NISAR L-band data as it becomes available from 2026;
(2) incorporating dynamic antecedent soil moisture and snowmelt proxies
as time-varying conditioning factors for seasonal susceptibility forecasting;
and (3) stratifying the model by trigger mechanism to separately characterise
monsoon-triggered, GLOF-triggered, and cloudburst-triggered susceptibility,
which exhibit different spatial patterns and require different mitigation strategies.
